% Two-panel figure for the polar (radial) model.
% Data written by models.pluto.jl -> figdata/  (col sep=comma).
% \input this from the standalone wrapper.  Swap the csv names / x=r->x=x
% and drop the second fields file for the two-defect (FFT) run.
% \footnotesize (8 pt at the 10pt base) is set on the tikzpicture so it cascades to
% all axis text -- matches the revtex4-2 APS caption size when the PDF is included at
% natural size (no rescaling). `font` is a TikZ key, not a pgfplots axis key.
%
% Measure the width of "-1.0" once into a length (pgfplots' `text width` chokes if you
% pass \widthof directly -- it runs the value through its math parser). \footnotesize
% here matches the tick-label font so the width is correct.
\ifdefined\ytickw\else\newlength{\ytickw}\fi
\settowidth{\ytickw}{\footnotesize$-1.0$}
% small-r analytic approximation (see derivation, "Behaviour near the core"):
%   rho = rhobar (1 + r^2/lam_rho^2),  p = r/lam_p,  v = r/tau_v
% constants (rhobar, 1/lam_p, 1/tau_v from the core fit; 1/lam_rho^2 from the density eq)
\def\rhobar{0.8149}\def\arho{0.2784}% \arho = 1/lambda_rho^2
\def\ap{1.2148}%                       \ap  = 1/lambda_p
\def\av{0.5416}%                       \av  = 1/tau_v
\begin{tikzpicture}[font=\footnotesize]
  \begin{groupplot}[
      % PRL single column = 8.6 cm wide
      group style={group size=1 by 2, vertical sep=1.1cm, group name=mainpanels},
      width=\prlcol, height=5cm,     % \prlcol = 8.6 cm; pgfplots `width` = total incl. labels
      xlabel=$r$,
      xmin=0, xmax=3,
      % reserve exactly the width of "-1.0" (measured above) for the y tick labels in
      % both panels, so the axis boxes -- and hence the two ylabels -- line up vertically
      yticklabel style={text width=\ytickw, align=right},
      legend style={draw=none, fill=none},
      legend cell align=left,
    ]
    % ---- panel 1: fields ----
    \nextgroupplot[ylabel={Fields}, ylabel style={name=ylaba}, ymin=0, legend pos=north west]
      \addplot[blue,  thick] table[x=r, y=rho]{figdata/polar_fields_rho.csv};
      \addlegendentry{$\rho$}
      \addplot[red,   thick] table[x=r, y=P]{figdata/polar_fields_Pv.csv};
      \addlegendentry{$p$}
      \addplot[black, thick] table[x=r, y=v]{figdata/polar_fields_Pv.csv};
      \addlegendentry{$v$}
      % small-r approximation (dashed, no legend entry)
      \addplot[blue,  thick, dashed, forget plot, domain=0:1.2, samples=40] {\rhobar*(1+\arho*x^2)};
      \addplot[red,   thick, dashed, forget plot, domain=0:1.0, samples=2]  {\ap*x};
      \addplot[black, thick, dashed, forget plot, domain=0:1.0, samples=2]  {\av*x};
    % ---- panel 2: density balance (rates) ----
    \nextgroupplot[ylabel={Density balance}, ylabel style={name=ylabb}, legend pos=north east]
      % blue = advection only = (total transport) - (diffusive part)
      \addplot[blue,  thick] table[x=r, y expr=\thisrow{transport}-\thisrow{diffusive}]{figdata/polar_balance.csv};
      \addlegendentry{$-\nabla\!\cdot(\boldsymbol v\rho)$}
      \addplot[red,   thick] table[x=r, y=diffusive]{figdata/polar_balance.csv};
      \addlegendentry{$M\Delta\mu$}
      \addplot[black, thick] table[x=r, y=turnover]{figdata/polar_balance.csv};
      \addlegendentry{$-\dfrac{\rho-\rho_0}{\tau}$}
  \end{groupplot}
  % panel labels, horizontally aligned with the ylabels (x from ylabel node, y from panel top)
  \node[anchor=south, yshift=-2mm] at (ylaba |- mainpanels c1r1.north) {(a)};
  \node[anchor=south, yshift=-2mm] at (ylabb |- mainpanels c1r2.north) {(b)};
  % --- inset in the fields panel: maximal velocity vs turnover time ---
  % data: figdata/polar_vsweep_summary.csv (from the tau-sweep in models.pluto.jl)
  \begin{axis}[
      at={(mainpanels c1r1.south east)}, anchor=south east, xshift=-3pt, yshift=10mm,
      width=3.8cm, height=2.8cm,
      xmin=0, ymin=0,
      xlabel=$\tau$, ylabel={$\max_r v$},
      xlabel style={font=\scriptsize, yshift=2mm}, ylabel style={font=\scriptsize},
      tick label style={font=\scriptsize},
      axis background/.style={fill=white},
    ]
    \addplot[black, thick] table[x=tau, y=vmax]{figdata/polar_vsweep_summary.csv};
  \end{axis}
\end{tikzpicture}
